% SEGMENT OLP-0375-S01
% Part: lambda-calculus
% Chapter: lambda-definablity
% Section: arithmetical-functions

\documentclass[../../../include/open-logic-section]{subfiles}

\begin{document}

\olfileid{lam}{rep}{arf}
\olsection{\usetoken{S}{lambda definable} எண்கணிதச் சார்புகள்}

\begin{prop}
  \ollabel{prop:succ-ld}
  தொடரிச் சார்பு~$\Succ$ ஆனது !!{lambda definable}.
\end{prop}

\begin{proof}
தொடரிச் சார்பை !!{lambda define}s ஒரு சொல்
\begin{align*}
  \fn{Succ} & \ident \lambd[a][\lambd[fx][f ({a} f x)]].
  \intertext{நமது குறுக்கக் குறிமான மரபுகளின்படி, இது பின்வருவதன் சுருக்கம்:}
  \fn{Succ} & \ident \lambd[a][\lambd[f][\lambd[x][(f (({a} f) x))]]].
  \intertext{$\fn{Succ}$ என்பது ஓர் எண்~$a$ ஐ உள்ளீடாக ஏற்று,
    $\lambd[fx][f ({a} f x)]$ என்ற மற்றொரு சார்பாக மதிப்பிடப்படும் ஒரு
    சார்பு. அந்தச் சார்பு தானாகவே ஒரு சர்ச் எண்ணுரு அல்ல. எனினும்,
    உள்ளீடு $a$ ஒரு சர்ச் எண்ணுருவாக இருந்தால், அது ஒரு சர்ச் எண்ணுருவாகக்
    குறைகிறது. பின்வருவதைக் கருதுக:}
   (\lambd[a][\lambd[fx][f ({a} f x)]])\,\num{n} & \redone
   \lambd[fx][f ({\num{n}} f x)].
   \intertext{உட்பொதிந்த சொல் $\num{n}fx$ ஒரு ரெடெக்ஸ்; ஏனெனில்
     $\num{n}$ என்பது $\lambd[fx][f^nx]$. ஆகவே
     $\num{n}fx \redone f^nx$; முழுச் சொல்லுக்கும்}
   \fn{Succ}\, \num n & \red \lambd[fx][f(f^n(x))],
\end{align*}
அதாவது, $\num{n+1}$ கிடைக்கிறது.
\end{proof}

\begin{ex}
  $\fn{Succ}$ ஐ~$\num{0}$ மீது, அதாவது $\lambd[fx][x]$ மீது,
  செயல்படுத்தும்போது என்ன நிகழ்கிறது என்று பார்ப்போம். சொற்களை முழுமையாக
  விரித்தெழுதுவோம்:
  \begin{align*}
    \fn{Succ}\, \num{0} & \ident (\lambd[a][\lambd[f][\lambd[x][(f (({a} f) x))]]])(\lambd[f][\lambd[x][x]])\\
    & \redone \lambd[f][\lambd[x][(f (({(\lambd[f][\lambd[x][x]])} f) x))]] \\
    & \redone \lambd[f][\lambd[x][(f ({(\lambd[x][x])} x))]] \\
    & \redone \lambd[f][\lambd[x][(f x)]] \ident \num{1}\\
  \end{align*}
\end{ex}

\begin{prob}
  சொல்
  \begin{align*}
    \fn{Succ}' & \ident \lambd[{n}][\lambd[fx][{n} f (f x)]]
  \end{align*}
  தொடரிச் சார்பை !!{lambda define}s. ஏன் என்று விளக்குக.
\end{prob}

\begin{prop}
  \ollabel{prop:add-ld}
  கூட்டல் சார்பு~$\Add$ ஆனது !!{lambda definable}.
\end{prop}

\begin{proof}
கூட்டலைப் பின்வரும் சொற்கள் !!{lambda define} செய்கின்றன:
\begin{align*}
  \fn{Add} & \ident \lambd[{a}{b}][\lambd[fx][{a} f ({b} f x)]]
\intertext{அல்லது, மாற்றாக,}
  \fn{Add}' & \ident \lambd[{a}{b}][{a}\, \fn{Succ} \, {b}].
\intertext{முதல் கூட்டல் பின்வருமாறு செயல்படுகிறது: $\fn{Add}$ முதலில்
  ${a}$, ${b}$ என்ற இரண்டு எண்களை ஏற்கிறது. அதன் விளைவு $f$, $x$
  ஆகியவற்றை ஏற்று $af(bfx)$ ஐத் திருப்பித்தரும் சார்பு. $a$, $b$ ஆகியவை
  $\num{n}$, $\num{m}$ என்ற சர்ச் எண்ணுருக்களாக இருந்தால், இது
  $f^{n+m}(x)$ ஆகக் குறைகிறது; இது $f^{n}(f^{m}(x))$ குச் சர்வசமம்.
  படிப்படியாக எழுதினால்:}
  (\lambd[{a}{b}][\lambd[fx][{a} f ({b} f x)]])\num n\,\num m & \redone
  \lambd[fx][\num{n}\, f (\num {m}\, f x)] \\
  & \redone \lambd[fx][\num{n}\, f (f^m x)] \\
  & \redone \lambd[fx][f^n (f^m x)] \ident \num {n+m}.
\intertext{கூட்டலின் இரண்டாம் பிரதிநிதித்துவமான $\fn{Add'}$ வேறுவிதமாகச்
  செயல்படுகிறது. $\num{n}$, $\num{m}$ என்ற இரண்டு சர்ச் எண்ணுருக்கள்மீது
  செயல்படுத்தினால்,}
\fn{Add}' \num n \,\num m
& \redone \num{n}\, \fn{Succ}\, \num{m}.
\intertext{ஆனால் $\num{n} f x$ எப்போதும் $f^n(x)$ ஆகக் குறைகிறது. ஆகவே,}
  \num{n}\,  \fn{Succ}\, \num{m} & \red \fn{Succ}^n(\num{m}).
\end{align*}
$\fn{Succ}$ தொடரிச் சார்பை !!{lambda define}s என்பதாலும், தொடரிச் சார்பை
$m$ மீது $n$ முறை செயல்படுத்துவது $n+m$ ஐத் தருவதாலும், இது இறுதியில்
~$\num{n+m}$ ஆகக் குறைகிறது.
\end{proof}

\begin{prop}
  \ollabel{prop:mult-ld}
  பெருக்கல் பின்வரும் சொல்லால் !!{lambda definable}:
  \[
  \fn{Mult} \ident \lambd[ab][\lambd[fx][a (b f) x]]
  \]
\end{prop}

\begin{proof}
  இது எவ்வாறு செயல்படுகிறது என்பதைக் காண, $\fn{Mult}$ ஐ
  $\num{n}$, $\num{m}$ என்ற சர்ச் எண்ணுருக்கள்மீது செயல்படுத்துவதாகக்
  கொள்க. $\fn{Mult} \, \num{n} \, \num{m}$ ஆனது
  $\lambd[fx][\num{n}(\num{m}\, f)x]$ ஆகக் குறைகிறது.
  சொல் $\num{m} f$ ஆனது தனது உள்ளீட்டின்மீது $f$ ஐ $m$ முறை
  செயல்படுத்தும் ஒரு சார்பை வரையறுக்கிறது. இதன் விளைவாக,
  $\num{n} (\num{m} f) x$ ஆனது “$f$ ஐ $m$ முறை செயல்படுத்து” என்ற
  சார்பையே~$x$ மீது $n$ முறை செயல்படுத்துகிறது. வேறு சொற்களில்,
  $f$ ஐ $x$ மீது $n\cdot m$ முறை செயல்படுத்துகிறோம். கிடைக்கும்
  இயல்புவடிவம் சர்ச் எண்ணுரு $\num{nm}$ மட்டுமே.
\end{proof}

\begin{editorial}
உண்மையில், $\eta$-குறைத்தலால் இச்சொல்லை மேலும் எளிமைப்படுத்தலாம்:
\[
  \fn{Mult} \ident \lambd[ab][\lambd[f][a (b f)]].
\]

ஆனால் அதற்கு முதலில் $\eta$-குறைத்தலை விளக்க வேண்டும்.
\end{editorial}

\begin{prob}
% SOURCE CORRECTION TA-LAM-035: iterate addition of b a times; the frozen term repeats a and therefore ignores b.
பின்வரும் சொல்லால் பெருக்கலை !!{lambda define} செய்யலாம்:
\[
  \fn{Mult}' \ident \lambd[ab][a (\fn{Add}\, b) \num{0}].
\]
இது ஏன் செயல்படுகிறது என்று விளக்குக.
\end{prob}

அடுக்கேற்றத்தின் $\lambd$-சொல் வரையறை வியக்கத்தக்க அளவுக்கு எளியது:
\[
  \fn{Exp} \ident \lambd[be][e b].
\]
முதல் உள்ளீடு $b$ அடிமானம்; இரண்டாம் உள்ளீடு~$e$ அடுக்கு. எண்களுக்கான நமது
குறியாக்கத்தின்படி $e f$ என்பது உள்ளுணர்வில் $f^e$. இதைப் புரிந்துகொள்வது
கடினமாக இருந்தால், மீள்செய்யப்பட்ட பெருக்கலால் அடுக்கேற்றத்தைப் பின்வருமாறும்
வரையறுக்கலாம்:
\[
  \fn{Exp}' \ident \lambd[be][e (\fn{Mult}\, b) \num{1}].
\]

சர்ச் எண்ணுருக்கள்மீதான முன்னி, கழித்தல் ஆகியவற்றின் வரையறைகள் நாம்
நினைப்பதுபோல் எளியவை அல்ல: அவற்றுக்கு ஜோடிகளின் குறியாக்கம் தேவை.

\end{document}
